% Part: normal-modal-logic
% Chapter: frame-definability
% Section: introduction

\documentclass[../../../include/open-logic-section]{subfiles}

\begin{document}

\olfileid{nml}{frd}{int}

\olsection{مقدمة}

من المسائل التي تهم المناطقة الموجهيين العلاقة بين علاقة الوصول وصدق
!!{formula}s معينة في النماذج ذات علاقة الوصول تلك. افترض، مثلًا، أن
علاقة الوصول انعكاسية، أي إنه لكل~$w \in W$ يكون~$Rww$. وبعبارة أخرى،
يمكن الوصول من كل عالم إلى نفسه. وهذا يعني أنه عندما تكون~$\Box !A$
صادقة في عالم~$w$، يكون~$w$ نفسه من العوالم التي يمكن الوصول إليها
والتي يجب أن تكون~$!A$ صادقة فيها. فإذا كانت علاقة الوصول~$R$ في
$\mModel{M}$ انعكاسية، فإن~$\Box !A \lif !A$ تكون صادقة أيًا كان
العالم~$w$ والصيغة~$!A$ اللذان نختارهما (وبعبارة أخرى، يكون
المخطط~$\Box p \lif p$ وجميع حالات إحلاله صادقة في~$\mModel{M}$).

لكن العكس كاذب. فليس صحيحًا، مثلًا، أنه إذا كانت~$\Box p \lif p$
صادقة في~$\mModel{M}$، كانت~$R$ انعكاسية. إذ يسهل أن نجد نموذجًا غير
انعكاسي~$\mModel{M}$ تكون فيه~$\Box p \lif p$ صادقة في جميع العوالم:
خذ نموذجًا ذا عالم وحيد~$w$ لا يمكن الوصول منه إلى نفسه، لكن
$w \in V(p)$. وباختيار قيمة صدق~$p$ على نحو مناسب، يمكننا جعل
$\Box p \lif p$ صادقة في نموذج غير انعكاسي.

والحل هو إخراج تعيين المتغيرات~$V$ من المعادلة. فإذا اشترطنا أن تكون
$\Box p \lif p$ صادقة في جميع عوالم~$\mModel{M}$، بصرف النظر عن
العوالم التي تنتمي إلى~$V(p)$، لزم أن تكون~$R$ انعكاسية. ففي أي نموذج
غير انعكاسي يوجد عالم واحد على الأقل~$w$ لا يصدق فيه~$Rww$. فإذا
جعلنا~$V(p) = W \setminus \{w\}$، كانت~$p$ صادقة في جميع العوالم غير
$w$، ومن ثم في جميع العوالم التي يمكن الوصول إليها من~$w$ (لأن
المضمون هو أنه لا يمكن الوصول من~$w$ إلى~$w$، و$w$ هو العالم الوحيد
الذي تكون فيه~$p$ كاذبة). ومن جهة أخرى، تكون~$p$ كاذبة في~$w$، فتكون
$\Box p \lif p$ كاذبة في~$w$.

يوحي هذا بأن نقدم رمزًا لبنى النماذج الخالية من التقييم؛ ونسميها
\emph{أطرًا}. فالإطار~$\mModel{F}$ ليس إلا زوجًا~$\tuple{W, R}$ يتألف
من مجموعة عوالم وعلاقة وصول. ونقول عندئذ إن كل نموذج
$\tuple{W, R, V}$ \emph{مبني على} الإطار~$\tuple{W, R}$. وبالعكس،
يحدد الإطار صنف النماذج المبنية عليه؛ ويحدد صنف من الأطر صنف النماذج
المبنية على أي إطار من ذلك الصنف. ويمكننا أن نعرّف
$\mModel{F} \Entails !A$، أي مفهوم كون !!a{formula} \emph{صالحة} في
إطار، كما يأتي:~$\mSat{M}{!A}$ لكل~$\mModel{M}$ مبني على~$\mModel{F}$.

وبهذا الترميز، يمكننا إقامة علاقات مطابقة بين ال!!{formula}s وأصناف
الأطر: فمثلًا،~$\mModel{F} \Entails \Box p \lif p$ إذا وفقط إذا كان
$\mModel{F}$ انعكاسيًا.

\end{document}
